% Generated. Do not edit.
% Projection scriptural-date-location/section, version 1: the Scriptural Date and
% Location page, generated whole from research/chronology.toml, the Scripture
% chronology corpus and the source library by tools/tpt document-types refresh.
% document: liturgy/roman-rite/1962/propers/temporal/61-twenty-first-after-pentecost
\runninghead{Scriptural Date and Location}
\section{Scriptural Date and Location}
\label{sec:date-location}

{\small Each passage the Mass reads or sings from Scripture, in canonical order. The Context column names what is dated --- the event the passage narrates, a setting it remembers, the traditional attribution of its author, or the composition of its book --- and where it belongs; the Date column gives the answer of the Scripture chronology corpus for that subject, with every alternative it records. Texts of the Mass that are not drawn from Scripture have no row. The evidence for each date follows the table.\par}
\medskip
% Generated. Do not edit.
% Deterministic projection of research/chronology.toml.
% schema: 5
% document: liturgy/roman-rite/1962/propers/temporal/61-twenty-first-after-pentecost
% generated-by: tools/tpt proper-chronology annotations
\newcommand{\chronologyannotationclaim}[6]{#6}
\newcommand{\chronologyannotationcomparisonclaim}[8]{#8}
\newcommand{\chronologyannotationreach}[3]{}
\newcommand{\chronologyannotationgroup}[3]{#3}
\newcommand{\chronologyannotationcomparisongroup}[4]{#4}
\newcommand{\chronologyannotation}[1]{%
  \ifcsname triptychchronologyannotation@#1\endcsname
    \csname triptychchronologyannotation@#1\endcsname
  \else
    \PackageError{triptych}{No chronology annotation for #1}{}%
  \fi}
% comparison-dependency: src/claude/liturgy/roman-rite/1962/propers/temporal/61-twenty-first-after-pentecost/research/chronology-profile-comparisons.toml sha256=1d7a538bc40cbe8f7d50f5136523fa3879fd042cf3965e244cdb9338c0d1f97b
% introit: Introit
\expandafter\def\csname triptychchronologyannotation@introit\endcsname{%
  \chronologyannotationgroup{historical-setting}{research-pending}{\textbf{Historical setting} -- Occasion date unresolved.}%
  \space \chronologyannotationgroup{composition}{nonuniform}{\textbf{Composition} -- No single date applies across every cited locus.}%
}
% epistle: Epistle
\expandafter\def\csname triptychchronologyannotation@epistle\endcsname{%
  \chronologyannotationgroup{composition}{disputed}{\textbf{Composition} -- disputed: \chronologyannotationreach{Eph.6.10}{Eph}{true}\chronologyannotationreach{Eph.6.11}{Eph}{true}\chronologyannotationreach{Eph.6.12}{Eph}{true}\chronologyannotationreach{Eph.6.13}{Eph}{true}\chronologyannotationreach{Eph.6.14}{Eph}{true}\chronologyannotationreach{Eph.6.15}{Eph}{true}\chronologyannotationreach{Eph.6.16}{Eph}{true}\chronologyannotationreach{Eph.6.17}{Eph}{true}\chronologyannotationclaim{composition.epistle-to-the-ephesians}{composition}{catholic-traditional-v1}{disputed}{a period between 58 and 63}{A.D. 58--63}; \chronologyannotationreach{Eph.6.10}{Eph}{true}\chronologyannotationreach{Eph.6.11}{Eph}{true}\chronologyannotationreach{Eph.6.12}{Eph}{true}\chronologyannotationreach{Eph.6.13}{Eph}{true}\chronologyannotationreach{Eph.6.14}{Eph}{true}\chronologyannotationreach{Eph.6.15}{Eph}{true}\chronologyannotationreach{Eph.6.16}{Eph}{true}\chronologyannotationreach{Eph.6.17}{Eph}{true}\chronologyannotationclaim{composition.epistle-to-the-ephesians}{composition}{catholic-traditional-v1}{disputed}{(Philemon; Colossians; Ephesians; Philippians), 61}{A.D. 61}.}%
  \space \chronologyannotationcomparisongroup{catholic-critical-v1}{composition}{composition-only}{\textbf{Composition (Catholic critical chronology, for comparison)} -- disputed: \chronologyannotationreach{Eph.6.10}{Eph}{true}\chronologyannotationreach{Eph.6.11}{Eph}{true}\chronologyannotationreach{Eph.6.12}{Eph}{true}\chronologyannotationreach{Eph.6.13}{Eph}{true}\chronologyannotationreach{Eph.6.14}{Eph}{true}\chronologyannotationreach{Eph.6.15}{Eph}{true}\chronologyannotationreach{Eph.6.16}{Eph}{true}\chronologyannotationreach{Eph.6.17}{Eph}{true}\chronologyannotationcomparisonclaim{critical.ephesians-later-disciple}{composition}{catholic-critical-v1}{catholic-critical-v1}{disputed}{passage.united-states-conference-of-catholic-bishops.new-american-bible-revised-edition.english-usccb-web-2026-09-28.ephesians-introduction}{around A.D. 80--100}{Ephesians under the later-disciple hypothesis (approximate), c. A.D. 80--100}.}%
}
% gradual: Gradual
\expandafter\def\csname triptychchronologyannotation@gradual\endcsname{%
  \chronologyannotationgroup{historical-setting}{research-pending}{\textbf{Historical setting} -- Occasion date unresolved.}%
  \space \chronologyannotationgroup{composition}{preferred}{\textbf{Composition}: \chronologyannotationreach{Ps.89.1}{Ps}{true}\chronologyannotationreach{Ps.89.2}{Ps}{true}\chronologyannotationclaim{critical.psalms.latest-composition-boundary}{composition}{catholic-critical-v1}{preferred}{before the Maccabean period, around 165 B.C.; no individual psalm can be dated securely}{Before c. 165 B.C.}}%
}
% alleluia: Alleluia
\expandafter\def\csname triptychchronologyannotation@alleluia\endcsname{%
  \chronologyannotationgroup{retrospective-event}{preferred}{\textbf{Retrospective event}: Preferred: \chronologyannotationreach{Ps.113.1}{Ps.77.12-13 Ps.77.42-53 Ps.104.26-38 Ps.104.43 Ps.105.7-12 Ps.105.21-22 Ps.113.1-5 Ps.135.10-15 Ps.80.6-7 2Esd.9.9-12 Acts.7.36}{false}\chronologyannotationclaim{israel.exodus.the-exodus}{retrospective-event}{catholic-traditional-v1}{preferred}{the four hundred and eightieth year after the children of Israel came out of the land of Egypt}{The Exodus from Egypt, the four hundred and eightieth year after the children of Israel came out of the land of Egypt}; alternative: \chronologyannotationreach{Ps.113.1}{Ps.77.12-13 Ps.77.42-53 Ps.104.26-38 Ps.104.43 Ps.105.7-12 Ps.105.21-22 Ps.113.1-5 Ps.135.10-15 Ps.80.6-7 2Esd.9.9-12 Acts.7.36}{false}\chronologyannotationclaim{israel.exodus.the-exodus}{retrospective-event}{catholic-traditional-v1}{alternate}{the Septuagint gives 440 years}{The Exodus from Egypt, the Septuagint gives 440 years}.}%
  \space \chronologyannotationgroup{composition}{preferred}{\textbf{Composition}: \chronologyannotationreach{Ps.113.1}{Ps}{true}\chronologyannotationclaim{critical.psalms.latest-composition-boundary}{composition}{catholic-critical-v1}{preferred}{before the Maccabean period, around 165 B.C.; no individual psalm can be dated securely}{Before c. 165 B.C.}}%
}
% gospel: Gospel
\expandafter\def\csname triptychchronologyannotation@gospel\endcsname{%
  \chronologyannotationgroup{narrated-event}{research-pending}{\textbf{Event} -- Narrated event date unresolved.}%
  \space \chronologyannotationgroup{composition}{disputed}{\textbf{Composition} -- disputed: \chronologyannotationreach{Matt.18.23}{Matt}{true}\chronologyannotationreach{Matt.18.24}{Matt}{true}\chronologyannotationreach{Matt.18.25}{Matt}{true}\chronologyannotationreach{Matt.18.26}{Matt}{true}\chronologyannotationreach{Matt.18.27}{Matt}{true}\chronologyannotationreach{Matt.18.28}{Matt}{true}\chronologyannotationreach{Matt.18.29}{Matt}{true}\chronologyannotationreach{Matt.18.30}{Matt}{true}\chronologyannotationreach{Matt.18.31}{Matt}{true}\chronologyannotationreach{Matt.18.32}{Matt}{true}\chronologyannotationreach{Matt.18.33}{Matt}{true}\chronologyannotationreach{Matt.18.34}{Matt}{true}\chronologyannotationreach{Matt.18.35}{Matt}{true}\chronologyannotationclaim{composition.gospel-of-matthew}{composition}{catholic-traditional-v1}{disputed}{about A.D. 38-45}{c. A.D. 38--45}; \chronologyannotationreach{Matt.18.23}{Matt}{true}\chronologyannotationreach{Matt.18.24}{Matt}{true}\chronologyannotationreach{Matt.18.25}{Matt}{true}\chronologyannotationreach{Matt.18.26}{Matt}{true}\chronologyannotationreach{Matt.18.27}{Matt}{true}\chronologyannotationreach{Matt.18.28}{Matt}{true}\chronologyannotationreach{Matt.18.29}{Matt}{true}\chronologyannotationreach{Matt.18.30}{Matt}{true}\chronologyannotationreach{Matt.18.31}{Matt}{true}\chronologyannotationreach{Matt.18.32}{Matt}{true}\chronologyannotationreach{Matt.18.33}{Matt}{true}\chronologyannotationreach{Matt.18.34}{Matt}{true}\chronologyannotationreach{Matt.18.35}{Matt}{true}\chronologyannotationclaim{composition.gospel-of-matthew}{composition}{catholic-traditional-v1}{disputed}{about the year 40-42}{c. A.D. 40--42}; \chronologyannotationreach{Matt.18.23}{Matt}{true}\chronologyannotationreach{Matt.18.24}{Matt}{true}\chronologyannotationreach{Matt.18.25}{Matt}{true}\chronologyannotationreach{Matt.18.26}{Matt}{true}\chronologyannotationreach{Matt.18.27}{Matt}{true}\chronologyannotationreach{Matt.18.28}{Matt}{true}\chronologyannotationreach{Matt.18.29}{Matt}{true}\chronologyannotationreach{Matt.18.30}{Matt}{true}\chronologyannotationreach{Matt.18.31}{Matt}{true}\chronologyannotationreach{Matt.18.32}{Matt}{true}\chronologyannotationreach{Matt.18.33}{Matt}{true}\chronologyannotationreach{Matt.18.34}{Matt}{true}\chronologyannotationreach{Matt.18.35}{Matt}{true}\chronologyannotationclaim{composition.gospel-of-matthew}{composition}{catholic-traditional-v1}{disputed}{the years 40-45}{A.D. 40--45}; \chronologyannotationreach{Matt.18.23}{Matt}{true}\chronologyannotationreach{Matt.18.24}{Matt}{true}\chronologyannotationreach{Matt.18.25}{Matt}{true}\chronologyannotationreach{Matt.18.26}{Matt}{true}\chronologyannotationreach{Matt.18.27}{Matt}{true}\chronologyannotationreach{Matt.18.28}{Matt}{true}\chronologyannotationreach{Matt.18.29}{Matt}{true}\chronologyannotationreach{Matt.18.30}{Matt}{true}\chronologyannotationreach{Matt.18.31}{Matt}{true}\chronologyannotationreach{Matt.18.32}{Matt}{true}\chronologyannotationreach{Matt.18.33}{Matt}{true}\chronologyannotationreach{Matt.18.34}{Matt}{true}\chronologyannotationreach{Matt.18.35}{Matt}{true}\chronologyannotationclaim{composition.gospel-of-matthew}{composition}{catholic-traditional-v1}{disputed}{about the year 60-68}{c. A.D. 60--68}; \chronologyannotationreach{Matt.18.23}{Matt}{true}\chronologyannotationreach{Matt.18.24}{Matt}{true}\chronologyannotationreach{Matt.18.25}{Matt}{true}\chronologyannotationreach{Matt.18.26}{Matt}{true}\chronologyannotationreach{Matt.18.27}{Matt}{true}\chronologyannotationreach{Matt.18.28}{Matt}{true}\chronologyannotationreach{Matt.18.29}{Matt}{true}\chronologyannotationreach{Matt.18.30}{Matt}{true}\chronologyannotationreach{Matt.18.31}{Matt}{true}\chronologyannotationreach{Matt.18.32}{Matt}{true}\chronologyannotationreach{Matt.18.33}{Matt}{true}\chronologyannotationreach{Matt.18.34}{Matt}{true}\chronologyannotationreach{Matt.18.35}{Matt}{true}\chronologyannotationclaim{composition.gospel-of-matthew}{composition}{catholic-traditional-v1}{disputed}{about the years 64-67}{c. A.D. 64--67}; \chronologyannotationreach{Matt.18.23}{Matt}{true}\chronologyannotationreach{Matt.18.24}{Matt}{true}\chronologyannotationreach{Matt.18.25}{Matt}{true}\chronologyannotationreach{Matt.18.26}{Matt}{true}\chronologyannotationreach{Matt.18.27}{Matt}{true}\chronologyannotationreach{Matt.18.28}{Matt}{true}\chronologyannotationreach{Matt.18.29}{Matt}{true}\chronologyannotationreach{Matt.18.30}{Matt}{true}\chronologyannotationreach{Matt.18.31}{Matt}{true}\chronologyannotationreach{Matt.18.32}{Matt}{true}\chronologyannotationreach{Matt.18.33}{Matt}{true}\chronologyannotationreach{Matt.18.34}{Matt}{true}\chronologyannotationreach{Matt.18.35}{Matt}{true}\chronologyannotationclaim{composition.gospel-of-matthew}{composition}{catholic-traditional-v1}{disputed}{about the year 50}{c. A.D. 50}.}%
  \space \chronologyannotationcomparisongroup{catholic-critical-v1}{composition}{composition-only}{\textbf{Composition (Catholic critical chronology, for comparison)}: \chronologyannotationreach{Matt.18.23}{Matt}{true}\chronologyannotationreach{Matt.18.24}{Matt}{true}\chronologyannotationreach{Matt.18.25}{Matt}{true}\chronologyannotationreach{Matt.18.26}{Matt}{true}\chronologyannotationreach{Matt.18.27}{Matt}{true}\chronologyannotationreach{Matt.18.28}{Matt}{true}\chronologyannotationreach{Matt.18.29}{Matt}{true}\chronologyannotationreach{Matt.18.30}{Matt}{true}\chronologyannotationreach{Matt.18.31}{Matt}{true}\chronologyannotationreach{Matt.18.32}{Matt}{true}\chronologyannotationreach{Matt.18.33}{Matt}{true}\chronologyannotationreach{Matt.18.34}{Matt}{true}\chronologyannotationreach{Matt.18.35}{Matt}{true}\chronologyannotationcomparisonclaim{critical.gospel-of-matthew}{composition}{catholic-critical-v1}{catholic-critical-v1}{preferred}{passage.united-states-conference-of-catholic-bishops.new-american-bible-revised-edition.english-usccb-web-2026-09-21.matthew-introduction}{post-A.D. 70 date}{The Greek Gospel of Matthew in the NABRE introduction, post-A.D. 70 date}.}%
}
% offertory: Offertory
\expandafter\def\csname triptychchronologyannotation@offertory\endcsname{%
  \chronologyannotationgroup{composition}{preferred}{\textbf{Composition}: \chronologyannotationreach{Job.1.1}{Job}{true}\chronologyannotationreach{Job.1.12}{Job}{true}\chronologyannotationreach{Job.2.6}{Job}{true}\chronologyannotationreach{Job.2.7}{Job}{true}\chronologyannotationclaim{critical.book-of-job}{composition}{catholic-critical-v1}{preferred}{probably sometime between the seventh and fifth centuries B.C.}{Probably sometime between the seventh and fifth centuries B.C.}}%
}
% communion: Communion
\expandafter\def\csname triptychchronologyannotation@communion\endcsname{%
  \chronologyannotationgroup{traditional-attribution}{disputed}{\textbf{Traditional attribution} -- disputed: \chronologyannotationreach{Ps.118.81}{Ps.118}{true}\chronologyannotationreach{Ps.118.84}{Ps.118}{true}\chronologyannotationreach{Ps.118.86}{Ps.118}{true}\chronologyannotationclaim{israel.monarchy.david-traditional-era}{traditional-attribution}{catholic-traditional-v1}{disputed}{reigned from 1055 to 1015 B.C.}{David (reign in the usual chronology), B.C. 1055--1015}.}%
  \space \chronologyannotationgroup{historical-setting}{research-pending}{\textbf{Historical setting} -- Occasion date unresolved.}%
  \space \chronologyannotationgroup{composition}{preferred}{\textbf{Composition}: \chronologyannotationreach{Ps.118.81}{Ps}{true}\chronologyannotationreach{Ps.118.84}{Ps}{true}\chronologyannotationreach{Ps.118.86}{Ps}{true}\chronologyannotationclaim{critical.psalms.latest-composition-boundary}{composition}{catholic-critical-v1}{preferred}{before the Maccabean period, around 165 B.C.; no individual psalm can be dated securely}{Before c. 165 B.C.}}%
}
% presentation: orientation-v3
\newcommand{\chronologyorientationrow}[2]{%
  \ifcsname triptychchronologyorientation@#1\endcsname
    \csname triptychchronologyorientation@#1\endcsname{#2}%
  \else
    \PackageError{triptych}{No chronology orientation row for #1}{}%
  \fi}
\expandafter\def\csname triptychchronologyorientation@row-2303cbdbb28a65b3\endcsname#1{%
Introit: Esth 13:9--11 & Composition: Esther; story at Susa; writing in Persia conjectured & \chronologyannotationreach{Esth.13.10}{Esth}{true}\chronologyannotationreach{Esth.13.11}{Esth}{true}\chronologyannotationreach{Esth.13.9}{Esth}{true}\chronologyannotationclaim{composition.book-of-esther}{composition}{catholic-traditional-v1}{preferred}{at the end of the reign of Xerxes I (485-465 B.C.) or during the reign of his son Artaxerxes I (465-425 B.C.)}{B.C.~465--425 (Conjectural early dating; later date possible)}\\[.3em]%
}
\expandafter\def\csname triptychchronologyorientation@row-b0e159e816e85ddc\endcsname#1{%
Offertory: Job 1:1, Job 1:12, Job 2:6--7 & Composition: Job; story in Hus; writing site unspecified & \chronologyannotationreach{Job.1.1}{Job}{true}\chronologyannotationreach{Job.1.12}{Job}{true}\chronologyannotationreach{Job.2.6}{Job}{true}\chronologyannotationreach{Job.2.7}{Job}{true}\chronologyannotationclaim{critical.book-of-job}{composition}{catholic-critical-v1}{preferred}{probably sometime between the seventh and fifth centuries B.C.}{Probably sometime between the seventh and fifth centuries B.C.}\\[.3em]%
}
\expandafter\def\csname triptychchronologyorientation@row-a873fa9c2a6cff33\endcsname#1{%
Gradual: Ps. 89:1--2 (Vg.; Heb. 90) & Historical setting & Occasion date unresolved\\*[.1em]%
}
\expandafter\def\csname triptychchronologyorientation@row-5c58d52409b5768a\endcsname#1{%
 & Composition: Psalms & \chronologyannotationreach{Ps.89.1}{Ps}{true}\chronologyannotationreach{Ps.89.2}{Ps}{true}\chronologyannotationclaim{critical.psalms.latest-composition-boundary}{composition}{catholic-critical-v1}{preferred}{before the Maccabean period, around 165 B.C.; no individual psalm can be dated securely}{Before c.~165 B.C.}\\[.3em]%
}
\expandafter\def\csname triptychchronologyorientation@row-ca449cd8277d8456\endcsname#1{%
Alleluia: Ps. 113:1 (Vg.; Heb. 114) & Retrospective event: The Exodus from Egypt; departure from Egypt; interval to Jerusalem's Temple & Preferred: \chronologyannotationreach{Ps.113.1}{Ps.77.12-13 Ps.77.42-53 Ps.104.26-38 Ps.104.43 Ps.105.7-12 Ps.105.21-22 Ps.113.1-5 Ps.135.10-15 Ps.80.6-7 2Esd.9.9-12 Acts.7.36}{false}\chronologyannotationclaim{israel.exodus.the-exodus}{retrospective-event}{catholic-traditional-v1}{preferred}{the four hundred and eightieth year after the children of Israel came out of the land of Egypt}{Temple begun in the 480th year after the Exodus (Vulgate)}; alternative: \chronologyannotationreach{Ps.113.1}{Ps.77.12-13 Ps.77.42-53 Ps.104.26-38 Ps.104.43 Ps.105.7-12 Ps.105.21-22 Ps.113.1-5 Ps.135.10-15 Ps.80.6-7 2Esd.9.9-12 Acts.7.36}{false}\chronologyannotationclaim{israel.exodus.the-exodus}{retrospective-event}{catholic-traditional-v1}{alternate}{the Septuagint gives 440 years}{Temple begun in the 440th year after the Exodus (Septuagint)}\\*[.1em]%
}
\expandafter\def\csname triptychchronologyorientation@row-f139a3b229cd1651\endcsname#1{%
 & Composition: Psalms & \chronologyannotationreach{Ps.113.1}{Ps}{true}\chronologyannotationclaim{critical.psalms.latest-composition-boundary}{composition}{catholic-critical-v1}{preferred}{before the Maccabean period, around 165 B.C.; no individual psalm can be dated securely}{Before c.~165 B.C.}\\[.3em]%
}
\expandafter\def\csname triptychchronologyorientation@row-642818c86ff9e033\endcsname#1{%
Introit: Ps 118:1 (Vg.; Heb. 119) & Traditional attribution: David's reign; Hebron, then Jerusalem & Disputed: \chronologyannotationreach{Ps.118.1}{Ps.118}{true}\chronologyannotationclaim{israel.monarchy.david-traditional-era}{traditional-attribution}{catholic-traditional-v1}{disputed}{reigned from 1055 to 1015 B.C.}{B.C.~1055--1015}\\*[.1em]%
}
\expandafter\def\csname triptychchronologyorientation@row-24a2902a47d0c589\endcsname#1{%
 & Historical setting & Occasion date unresolved\\*[.1em]%
}
\expandafter\def\csname triptychchronologyorientation@row-2c2dfef7006f9a46\endcsname#1{%
 & Composition: Psalms & \chronologyannotationreach{Ps.118.1}{Ps}{true}\chronologyannotationclaim{critical.psalms.latest-composition-boundary}{composition}{catholic-critical-v1}{preferred}{before the Maccabean period, around 165 B.C.; no individual psalm can be dated securely}{Before c.~165 B.C.}\\[.3em]%
}
\expandafter\def\csname triptychchronologyorientation@row-8fef59fec321f55b\endcsname#1{%
Communion: Ps. 118:81, 84, 86 (Vg.; Heb. 119) & Traditional attribution: David's reign; Hebron, then Jerusalem & Disputed: \chronologyannotationreach{Ps.118.81}{Ps.118}{true}\chronologyannotationreach{Ps.118.84}{Ps.118}{true}\chronologyannotationreach{Ps.118.86}{Ps.118}{true}\chronologyannotationclaim{israel.monarchy.david-traditional-era}{traditional-attribution}{catholic-traditional-v1}{disputed}{reigned from 1055 to 1015 B.C.}{B.C.~1055--1015}\\*[.1em]%
}
\expandafter\def\csname triptychchronologyorientation@row-020e5c17dd2119d2\endcsname#1{%
 & Historical setting & Occasion date unresolved\\*[.1em]%
}
\expandafter\def\csname triptychchronologyorientation@row-9aa748092505603b\endcsname#1{%
 & Composition: Psalms & \chronologyannotationreach{Ps.118.81}{Ps}{true}\chronologyannotationreach{Ps.118.84}{Ps}{true}\chronologyannotationreach{Ps.118.86}{Ps}{true}\chronologyannotationclaim{critical.psalms.latest-composition-boundary}{composition}{catholic-critical-v1}{preferred}{before the Maccabean period, around 165 B.C.; no individual psalm can be dated securely}{Before c.~165 B.C.}\\[.3em]%
}
\expandafter\def\csname triptychchronologyorientation@row-c83c455517aa2f92\endcsname#1{%
Gospel: Matt. 18:23--35 & Narrated event: Jesus tells the parable of the unforgiving servant; Capharnaum by narrative continuity & Date unresolved\\*[.1em]%
}
\expandafter\def\csname triptychchronologyorientation@row-29efed676e57212f\endcsname#1{%
 & Composition: Matthew; Semitic original associated with Palestine & Disputed: \chronologyannotationreach{Matt.18.23}{Matt}{true}\chronologyannotationreach{Matt.18.24}{Matt}{true}\chronologyannotationreach{Matt.18.25}{Matt}{true}\chronologyannotationreach{Matt.18.26}{Matt}{true}\chronologyannotationreach{Matt.18.27}{Matt}{true}\chronologyannotationreach{Matt.18.28}{Matt}{true}\chronologyannotationreach{Matt.18.29}{Matt}{true}\chronologyannotationreach{Matt.18.30}{Matt}{true}\chronologyannotationreach{Matt.18.31}{Matt}{true}\chronologyannotationreach{Matt.18.32}{Matt}{true}\chronologyannotationreach{Matt.18.33}{Matt}{true}\chronologyannotationreach{Matt.18.34}{Matt}{true}\chronologyannotationreach{Matt.18.35}{Matt}{true}\chronologyannotationclaim{composition.gospel-of-matthew}{composition}{catholic-traditional-v1}{disputed}{about A.D. 38-45}{c.~A.D.~38--45}; \chronologyannotationreach{Matt.18.23}{Matt}{true}\chronologyannotationreach{Matt.18.24}{Matt}{true}\chronologyannotationreach{Matt.18.25}{Matt}{true}\chronologyannotationreach{Matt.18.26}{Matt}{true}\chronologyannotationreach{Matt.18.27}{Matt}{true}\chronologyannotationreach{Matt.18.28}{Matt}{true}\chronologyannotationreach{Matt.18.29}{Matt}{true}\chronologyannotationreach{Matt.18.30}{Matt}{true}\chronologyannotationreach{Matt.18.31}{Matt}{true}\chronologyannotationreach{Matt.18.32}{Matt}{true}\chronologyannotationreach{Matt.18.33}{Matt}{true}\chronologyannotationreach{Matt.18.34}{Matt}{true}\chronologyannotationreach{Matt.18.35}{Matt}{true}\chronologyannotationclaim{composition.gospel-of-matthew}{composition}{catholic-traditional-v1}{disputed}{about the year 40-42}{c.~A.D.~40--42}; \chronologyannotationreach{Matt.18.23}{Matt}{true}\chronologyannotationreach{Matt.18.24}{Matt}{true}\chronologyannotationreach{Matt.18.25}{Matt}{true}\chronologyannotationreach{Matt.18.26}{Matt}{true}\chronologyannotationreach{Matt.18.27}{Matt}{true}\chronologyannotationreach{Matt.18.28}{Matt}{true}\chronologyannotationreach{Matt.18.29}{Matt}{true}\chronologyannotationreach{Matt.18.30}{Matt}{true}\chronologyannotationreach{Matt.18.31}{Matt}{true}\chronologyannotationreach{Matt.18.32}{Matt}{true}\chronologyannotationreach{Matt.18.33}{Matt}{true}\chronologyannotationreach{Matt.18.34}{Matt}{true}\chronologyannotationreach{Matt.18.35}{Matt}{true}\chronologyannotationclaim{composition.gospel-of-matthew}{composition}{catholic-traditional-v1}{disputed}{the years 40-45}{A.D.~40--45}; \chronologyannotationreach{Matt.18.23}{Matt}{true}\chronologyannotationreach{Matt.18.24}{Matt}{true}\chronologyannotationreach{Matt.18.25}{Matt}{true}\chronologyannotationreach{Matt.18.26}{Matt}{true}\chronologyannotationreach{Matt.18.27}{Matt}{true}\chronologyannotationreach{Matt.18.28}{Matt}{true}\chronologyannotationreach{Matt.18.29}{Matt}{true}\chronologyannotationreach{Matt.18.30}{Matt}{true}\chronologyannotationreach{Matt.18.31}{Matt}{true}\chronologyannotationreach{Matt.18.32}{Matt}{true}\chronologyannotationreach{Matt.18.33}{Matt}{true}\chronologyannotationreach{Matt.18.34}{Matt}{true}\chronologyannotationreach{Matt.18.35}{Matt}{true}\chronologyannotationclaim{composition.gospel-of-matthew}{composition}{catholic-traditional-v1}{disputed}{about the year 60-68}{c.~A.D.~60--68}; \chronologyannotationreach{Matt.18.23}{Matt}{true}\chronologyannotationreach{Matt.18.24}{Matt}{true}\chronologyannotationreach{Matt.18.25}{Matt}{true}\chronologyannotationreach{Matt.18.26}{Matt}{true}\chronologyannotationreach{Matt.18.27}{Matt}{true}\chronologyannotationreach{Matt.18.28}{Matt}{true}\chronologyannotationreach{Matt.18.29}{Matt}{true}\chronologyannotationreach{Matt.18.30}{Matt}{true}\chronologyannotationreach{Matt.18.31}{Matt}{true}\chronologyannotationreach{Matt.18.32}{Matt}{true}\chronologyannotationreach{Matt.18.33}{Matt}{true}\chronologyannotationreach{Matt.18.34}{Matt}{true}\chronologyannotationreach{Matt.18.35}{Matt}{true}\chronologyannotationclaim{composition.gospel-of-matthew}{composition}{catholic-traditional-v1}{disputed}{about the years 64-67}{c.~A.D.~64--67}; \chronologyannotationreach{Matt.18.23}{Matt}{true}\chronologyannotationreach{Matt.18.24}{Matt}{true}\chronologyannotationreach{Matt.18.25}{Matt}{true}\chronologyannotationreach{Matt.18.26}{Matt}{true}\chronologyannotationreach{Matt.18.27}{Matt}{true}\chronologyannotationreach{Matt.18.28}{Matt}{true}\chronologyannotationreach{Matt.18.29}{Matt}{true}\chronologyannotationreach{Matt.18.30}{Matt}{true}\chronologyannotationreach{Matt.18.31}{Matt}{true}\chronologyannotationreach{Matt.18.32}{Matt}{true}\chronologyannotationreach{Matt.18.33}{Matt}{true}\chronologyannotationreach{Matt.18.34}{Matt}{true}\chronologyannotationreach{Matt.18.35}{Matt}{true}\chronologyannotationclaim{composition.gospel-of-matthew}{composition}{catholic-traditional-v1}{disputed}{about the year 50}{c.~A.D.~50}\\*[.1em]%
}
\expandafter\def\csname triptychchronologyorientation@row-751d91210128bba1\endcsname#1{%
 & Composition, Catholic critical comparison: Greek Matthew; Antioch suggested & \chronologyannotationreach{Matt.18.23}{Matt}{true}\chronologyannotationreach{Matt.18.24}{Matt}{true}\chronologyannotationreach{Matt.18.25}{Matt}{true}\chronologyannotationreach{Matt.18.26}{Matt}{true}\chronologyannotationreach{Matt.18.27}{Matt}{true}\chronologyannotationreach{Matt.18.28}{Matt}{true}\chronologyannotationreach{Matt.18.29}{Matt}{true}\chronologyannotationreach{Matt.18.30}{Matt}{true}\chronologyannotationreach{Matt.18.31}{Matt}{true}\chronologyannotationreach{Matt.18.32}{Matt}{true}\chronologyannotationreach{Matt.18.33}{Matt}{true}\chronologyannotationreach{Matt.18.34}{Matt}{true}\chronologyannotationreach{Matt.18.35}{Matt}{true}\chronologyannotationcomparisonclaim{critical.gospel-of-matthew}{composition}{catholic-critical-v1}{catholic-critical-v1}{preferred}{passage.united-states-conference-of-catholic-bishops.new-american-bible-revised-edition.english-usccb-web-2026-09-21.matthew-introduction}{post-A.D. 70 date}{Post-A.D.~70 date}\\[.3em]%
}
\expandafter\def\csname triptychchronologyorientation@row-c6dbeefe691eae1f\endcsname#1{%
Epistle: Eph. 6:10--17 & Composition: Ephesians; Rome or Caesarea; recipients in Asia Minor & Disputed: \chronologyannotationreach{Eph.6.10}{Eph}{true}\chronologyannotationreach{Eph.6.11}{Eph}{true}\chronologyannotationreach{Eph.6.12}{Eph}{true}\chronologyannotationreach{Eph.6.13}{Eph}{true}\chronologyannotationreach{Eph.6.14}{Eph}{true}\chronologyannotationreach{Eph.6.15}{Eph}{true}\chronologyannotationreach{Eph.6.16}{Eph}{true}\chronologyannotationreach{Eph.6.17}{Eph}{true}\chronologyannotationclaim{composition.epistle-to-the-ephesians}{composition}{catholic-traditional-v1}{disputed}{a period between 58 and 63}{A.D.~58--63}; \chronologyannotationreach{Eph.6.10}{Eph}{true}\chronologyannotationreach{Eph.6.11}{Eph}{true}\chronologyannotationreach{Eph.6.12}{Eph}{true}\chronologyannotationreach{Eph.6.13}{Eph}{true}\chronologyannotationreach{Eph.6.14}{Eph}{true}\chronologyannotationreach{Eph.6.15}{Eph}{true}\chronologyannotationreach{Eph.6.16}{Eph}{true}\chronologyannotationreach{Eph.6.17}{Eph}{true}\chronologyannotationclaim{composition.epistle-to-the-ephesians}{composition}{catholic-traditional-v1}{disputed}{(Philemon; Colossians; Ephesians; Philippians), 61}{A.D.~61}\\*[.1em]%
}
\expandafter\def\csname triptychchronologyorientation@row-523e148781ef9e94\endcsname#1{%
 & Composition, Catholic critical comparison: Ephesians (later-disciple hypothesis) & Disputed: \chronologyannotationreach{Eph.6.10}{Eph}{true}\chronologyannotationreach{Eph.6.11}{Eph}{true}\chronologyannotationreach{Eph.6.12}{Eph}{true}\chronologyannotationreach{Eph.6.13}{Eph}{true}\chronologyannotationreach{Eph.6.14}{Eph}{true}\chronologyannotationreach{Eph.6.15}{Eph}{true}\chronologyannotationreach{Eph.6.16}{Eph}{true}\chronologyannotationreach{Eph.6.17}{Eph}{true}\chronologyannotationcomparisonclaim{critical.ephesians-later-disciple}{composition}{catholic-critical-v1}{catholic-critical-v1}{disputed}{passage.united-states-conference-of-catholic-bishops.new-american-bible-revised-edition.english-usccb-web-2026-09-28.ephesians-introduction}{around A.D. 80--100}{c.~A.D.~80--100}\\[.3em]%
}

% Generated. Do not edit.
% Complete orientation-v3 table; wording, dates and order belong to the corpus.
\begin{chronologysharedtable}
\chronologyorientationrow{row-2303cbdbb28a65b3}{}
\midrule
\chronologyorientationrow{row-b0e159e816e85ddc}{}
\midrule
\chronologyorientationrow{row-a873fa9c2a6cff33}{}
\chronologyorientationrow{row-5c58d52409b5768a}{}
\midrule
\chronologyorientationrow{row-ca449cd8277d8456}{}
\chronologyorientationrow{row-f139a3b229cd1651}{}
\midrule
\chronologyorientationrow{row-3ae21b72a8dbc7d2}{}
\chronologyorientationrow{row-54be670c8df1af30}{}
\chronologyorientationrow{row-a69079f4678c4dc5}{}
\midrule
\chronologyorientationrow{row-4f57d8133d9298cc}{}
\chronologyorientationrow{row-a5d51e9ca9b040bf}{}
\chronologyorientationrow{row-07c9bb9d0eca0d6f}{}
\midrule
\chronologyorientationrow{row-c83c455517aa2f92}{}
\chronologyorientationrow{row-29efed676e57212f}{}
\chronologyorientationrow{row-751d91210128bba1}{}
\midrule
\chronologyorientationrow{row-c6dbeefe691eae1f}{}
\chronologyorientationrow{row-523e148781ef9e94}{}
\end{chronologysharedtable}
{\footnotesize Psalm numbering: Vg. = Vulgate; Heb. = Hebrew.\par}


\subsection{Evidence for each date}
\label{sec:date-location-evidence}
{\small For each passage, as the Scriptural Date and Location table names it: what is dated, how the corpus weighs each answer, the source's own words, and the sources cited. A subject the corpus does not date is reported as such and given no date.\par}
{\small
\par\smallskip\noindent\textit{Introit: Esth. 13:9--11}\par\nobreak
\begin{itemize}[leftmargin=1.2em,itemsep=1pt,topsep=1pt,parsep=0pt]
\item \textbf{Composition} (preferred). The Book of Esther: B.C.~465--425 (Conjectural early dating; later date possible) (``at the end of the reign of Xerxes I (485-465 B.C.) or during the reign of his son Artaxerxes I (465-425 B.C.)''; The Catholic Encyclopedia, Volume 5 (New York, 1909)). Place: story at Susa; writing in Persia conjectured (The Catholic Encyclopedia, Volume 5 (New York, 1909); The Holy Bible, Douay-Rheims (Challoner revision), Project Gutenberg eBook 1581: Esth. 1:2).
\end{itemize}
\par\smallskip\noindent\textit{Offertory: Job 1:1, Job 1:12, Job 2:6--7}\par\nobreak
\begin{itemize}[leftmargin=1.2em,itemsep=1pt,topsep=1pt,parsep=0pt]
\item \textbf{Composition} (preferred). The Book of Job: Probably sometime between the seventh and fifth centuries B.C. (Catholic critical chronology; New American Bible, Revised Edition --- USCCB web state (2026), Job introduction). Place: story in Hus; writing site unspecified (The Holy Bible, Douay-Rheims (Challoner revision), Project Gutenberg eBook 1581: Job 1:1).
\end{itemize}
\par\smallskip\noindent\textit{Gradual: Ps. 89:1--2 (Vg.; Heb. 90)}\par\nobreak
\begin{itemize}[leftmargin=1.2em,itemsep=1pt,topsep=1pt,parsep=0pt]
\item \textbf{Historical setting}. No date: the chronology corpus has not yet dated this (research pending).
\item \textbf{Composition} (preferred). Psalms: Before c.~165 B.C. (``before the Maccabean period, around 165 B.C.; no individual psalm can be dated securely''; Catholic critical chronology; New American Bible, Revised Edition --- USCCB web state (2026), Psalms introduction).
\end{itemize}
\par\smallskip\noindent\textit{Alleluia: Ps. 113:1 (Vg.; Heb. 114)}\par\nobreak
\begin{itemize}[leftmargin=1.2em,itemsep=1pt,topsep=1pt,parsep=0pt]
\item \textbf{Retrospective event} (answers differ). The Exodus from Egypt: preferred: Temple begun in the 480th year after the Exodus (Vulgate) (``the four hundred and eightieth year after the children of Israel came out of the land of Egypt''; The Holy Bible, Douay-Rheims (Challoner revision), Project Gutenberg eBook 1581: 3 Kings 6:1); alternative: Temple begun in the 440th year after the Exodus (Septuagint) (``the Septuagint gives 440 years''; The Catholic Encyclopedia, Volume 3 (New York, 1908)). Place: departure from Egypt; interval to Jerusalem's Temple (The Holy Bible, Douay-Rheims (Challoner revision), Project Gutenberg eBook 1581: Ps. 113:1, 3 Kings 6:1, 2 Par. 3:1).
\item \textbf{Composition}. As under Gradual: Ps. 89:1--2 (Vg.; Heb. 90).
\end{itemize}
\par\smallskip\noindent\textit{Introit: Ps. 118:1 (Vg.; Heb. 119)}\par\nobreak
\begin{itemize}[leftmargin=1.2em,itemsep=1pt,topsep=1pt,parsep=0pt]
\item \textbf{Traditional attribution} (disputed). David's reign in the usual chronology: B.C.~1055--1015 (Corbett also reports writers dating the reign 30--50 years later) (``reigned from 1055 to 1015 B.C.''; The Catholic Encyclopedia, Volume 4 (New York, 1908), article 'David, King,' printed p. 642, opening paragraph). Place: Hebron, then Jerusalem (The Holy Bible, Douay-Rheims (Challoner revision), Project Gutenberg eBook 1581: 2 Kings 5:5).
\item \textbf{Historical setting}. As under Gradual: Ps. 89:1--2 (Vg.; Heb. 90).
\item \textbf{Composition}. As under Gradual: Ps. 89:1--2 (Vg.; Heb. 90).
\end{itemize}
\par\smallskip\noindent\textit{Communion: Ps. 118:81, 84, 86 (Vg.; Heb. 119)}\par\nobreak
\begin{itemize}[leftmargin=1.2em,itemsep=1pt,topsep=1pt,parsep=0pt]
\item \textbf{Traditional attribution}. As under Introit: Ps. 118:1 (Vg.; Heb. 119).
\item \textbf{Historical setting}. As under Gradual: Ps. 89:1--2 (Vg.; Heb. 90).
\item \textbf{Composition}. As under Gradual: Ps. 89:1--2 (Vg.; Heb. 90).
\end{itemize}
\par\smallskip\noindent\textit{Gospel: Matt. 18:23--35}\par\nobreak
\begin{itemize}[leftmargin=1.2em,itemsep=1pt,topsep=1pt,parsep=0pt]
\item \textbf{Narrated event}. Jesus tells the parable of the unforgiving servant; Capharnaum by narrative continuity (The Holy Bible, Douay-Rheims (Challoner revision), Project Gutenberg eBook 1581: Matt. 17:23, 18:1, 19:1, 18:21, 18:22, 18:23, 18:35). No date: the chronology corpus has not yet dated this (research pending).
\item \textbf{Composition} (disputed). The Gospel of St Matthew: c.~A.D.~38--45 (``about A.D. 38-45''; The Catholic Encyclopedia, Volume 10 (1911); The Catholic Encyclopedia, Volume 14 (New York, 1912)); c.~A.D.~40--42 (``about the year 40-42''; The Catholic Encyclopedia, Volume 10 (1911)); A.D.~40--45 (``the years 40-45''; The Catholic Encyclopedia, Volume 10 (1911)); c.~A.D.~60--68 (``about the year 60-68''; The Catholic Encyclopedia, Volume 10 (1911)); c.~A.D.~64--67 (``about the years 64-67''; The Catholic Encyclopedia, Volume 10 (1911)); c.~A.D.~50 (``about the year 50''; The Catholic Encyclopedia, Volume 14 (New York, 1912)). Place: Semitic original associated with Palestine (The Catholic Encyclopedia, Volume 10 (1911)).
\item \textbf{Composition, Catholic critical comparison} (preferred). The Greek Gospel of Matthew in the NABRE introduction: Post-A.D.~70 date (Catholic critical chronology; New American Bible, Revised Edition: USCCB web state (2026), Matthew introduction, final five prose paragraphs on authorship, sources, composition and place; narrative overview of chapters 19-25). Place: Antioch suggested (New American Bible, Revised Edition: USCCB web state (2026), Matthew introduction, final five prose paragraphs on authorship, sources, composition and place; narrative overview of chapters 19-25).
\end{itemize}
\par\smallskip\noindent\textit{Epistle: Eph. 6:10--17}\par\nobreak
\begin{itemize}[leftmargin=1.2em,itemsep=1pt,topsep=1pt,parsep=0pt]
\item \textbf{Composition} (disputed). The Epistle to the Ephesians: A.D.~58--63 (``a period between 58 and 63''; The Catholic Encyclopedia, Volume 5 (New York, 1909)); A.D.~61 (``(Philemon; Colossians; Ephesians; Philippians), 61''; The Catholic Encyclopedia, Volume 11 (New York, 1911)). Place: Rome or Caesarea; recipients in Asia Minor (The Catholic Encyclopedia, Volume 5 (New York, 1909)).
\item \textbf{Composition, Catholic critical comparison} (disputed). Ephesians under the later-disciple hypothesis (approximate): c.~A.D.~80--100 (``around A.D. 80--100''; Catholic critical chronology; New American Bible, Revised Edition: USCCB web state (2026), The Letter to the Ephesians, complete introduction; especially prose paragraphs 3-4, physical HTML lines 626 and 628).
\end{itemize}
\par}
